\section{Introduction}
\label{sec:intro}

Every decade following the census, state legislatures draw three sets of
district boundaries for the same geographic territory: congressional districts
(determined by federal apportionment), state senate districts, and state house
districts.  Under current practice these three processes are legally and
computationally independent.  A legislature controlling all three maps can
exploit this independence to reinforce partisan advantages across chambers: pack
a geographic community into a safe district at the congressional level, pack the
same community into a safe senate district, and pack it again at the house
level.  No single anti-gerrymandering doctrine spans all three chambers
simultaneously.

\subsection{The Structural Independence Problem}

Consider North Carolina after 2020 apportionment: $C = 14$ congressional
seats, $S = 50$ senate seats, $H = 120$ house seats.  A partisan legislature
facing a competitive map at one level can draw the other two levels to
compensate.  If a court strikes down the congressional map, the partisan sorting
of communities is preserved in the senate and house maps, which then serve as
the foundation for the next round of congressional redistricting.  The maps are
legally independent but geographically entangled.

The \emph{structural} version of this problem is even more basic.  The three
seat counts $(C, S, H)$ rarely share a compatible prime factorization.  There is
in general no way to recurse a 14-seat bisection tree to produce exactly 50
senate children and 120 house grandchildren.  Each chamber's map is drawn by a
different algorithm, consultant, or commission, with different geographic
anchors, producing three maps that are at best loosely correlated.

\subsection{NestSection}

We propose \emph{NestSection}, an algorithm that forces the three chambers to
share a common \emph{compatible factorization spine}.  The spine is the
coarsest geographic division that all three chambers agree on: it is
determined by $g = \gcd(C, S, H)$, the greatest common divisor of the three
seat counts.  Within each of the $g$ spine regions, the three chambers recurse
independently, but they share the same geographic boundaries at the top level.

The key insight is that $g$ determines how much cross-chamber structure can
be shared.  When $g = \min(C, S, H)$, the smallest chamber's seat count divides
both larger counts exactly, and the nesting is perfect: every district of the
smaller chamber is a union of districts from the larger chambers.  When $g = 1$,
no non-trivial common structure exists and the three chambers can only share a
root node (the state boundary itself).

\paragraph{Formal compatibility score.}  We define the
\emph{spine compatibility score} as:
\begin{equation}
  \label{eq:score}
  \sigma(C, S, H) \;=\; \left(1 - \frac{\gcd(C,S,H)}{\min(C,S,H)}\right) \times 100 \;\in\; [0, 100].
\end{equation}
A score of 0 means the smallest chamber is entirely captured by the shared
trunk; a score of 100 means the trunk is trivially a single root.

\paragraph{Main contributions.}
\begin{enumerate}
  \item A formal definition of compatible factorization spines and a
        polynomial-time algorithm ($\spine$) for their construction
        (Section~\ref{sec:algorithm}).
  \item A 50-state compatibility census under the 2020 apportionment, including
        an exact classification of which states admit strict nesting and which
        require best-effort Mode~3 nesting (Section~\ref{sec:evaluation}).
  \item A legal argument that shared-spine redistricting is an anti-gerrymandering
        mechanism by construction, and a statutory design proposal for states
        where $g \geq 5$ (Section~\ref{sec:discussion}).
\end{enumerate}

\paragraph{Roadmap.}  Section~\ref{sec:related} reviews multi-member district
theory, nested districting proposals, and the gerrymandering literature
relevant to cross-chamber coherence.  Section~\ref{sec:algorithm} defines
the $\spine$ algorithm formally, proves its key properties, and describes the
three nesting modes.  Section~\ref{sec:evaluation} presents the 50-state
compatibility census and three case studies (Oregon, North Carolina, Texas).
Section~\ref{sec:discussion} develops the legal argument, addresses limitations,
and outlines future work.  Section~\ref{sec:conclusion} summarizes.

\paragraph{Implementation note.}  The $\spine$ algorithm is implemented in the
\texttt{redist-apportion} crate of the \texttt{redist} Rust workspace; the
\texttt{compatible\_spines}, \texttt{spine\_compatibility\_score}, and
\texttt{us\_state\_compatibility\_table} functions are publicly exported and
tested against all 50 states.
